% !TeX root = ../main.tex
\section{Comparison, Failure Modes and Open Problems}
\label{sec:comparison}

\subsection{Comparison}
Table~\ref{tab:methods} places the main methods by input, output and
availability. The pattern is a split: methods that handle real mixes
(MuScriptor, YourMT3+, source separators) output notes or stems without
fingering, and methods that output tablature (TabCNN, FretNet, Noise2Fret,
TART) require solo guitar. A modular pipeline that chains a separator, a
guitar note model and a fingering model is currently the only way to obtain
tablature from a commercial recording, and its accuracy is bounded by the
weakest stage and by error accumulation, as TART's drop from 71.8\% to
54.08\% Tab F1 illustrates.

\begin{table*}[t]
\centering
\caption{Main methods by stage. ``Mix'': designed or shown to work on
multi-instrument recordings. Reported results use different test sets and are
not directly comparable.}
\label{tab:methods}
\footnotesize
\setlength{\tabcolsep}{3.5pt}
\begin{tabular}{@{}llllll@{}}
\toprule
Method & Stage & Input & Output & Mix & Reported (as published) \\
\midrule
HT Demucs~\cite{rouard2023htdemucs} & separation & mix & 4 or 6 stems (guitar in 6s) & yes & 9.20\,dB SDR, MUSDB18 \\
BS-RoFormer~\cite{lu2023bsroformer} & separation & mix & stems & yes & 9.80\,dB SDR, MUSDB18HQ; SDX23 winner \\
Banquet~\cite{watcharasupat2024banquet} & separation & mix & queried stem & yes & beats htdemucs\_6s on guitar, MoisesDB \\
Kong et al.~\cite{kong2021highres} & notes & solo piano & notes & no & 96.72\% onset F1, MAESTRO \\
GAPS model~\cite{riley2024gaps} & notes & solo guitar & notes & no & SOTA on GuitarSet, zero-shot \\
TIT~\cite{chen2025egdbpg} & notes & electric guitar & notes & no & gains across amp types \\
Basic Pitch~\cite{bittner2022basicpitch} & notes & any instrument & notes, contour & no & lightweight baseline \\
MT3~\cite{gardner2022mt3} & notes & mix & notes + instrument & yes & large gains on guitar \\
YourMT3+~\cite{chang2024yourmt3} & notes & mix & notes + instrument & yes & 10 datasets; limited on pop \\
MuScriptor~\cite{rouard2026muscriptor} & notes & mix & notes + instrument & yes & 11k\,h real data; 60.4 onset F1 (own test) \\
Beat This!~\cite{foscarin2024beatthis} & rhythm & mix or solo & beats, downbeats & yes & SOTA beat F1 without DBN \\
AlignBeat~\cite{ahn2025alignbeat} & rhythm & mix & beats, downbeats & yes & +3.6 downbeat CMLt over Beat This! \\
Viterbi/HMM~\cite{sayegh1989fingering,hori2016minimax} & fingering & notes & string, fret & -- & hand-designed costs \\
MIDI-to-Tab~\cite{edwards2024miditotab} & fingering & notes & string & -- & preferred in user study \\
Fretting-Transformer~\cite{hamberger2025fretting} & fingering & notes & string, fret & -- & beats A* and Guitar Pro \\
Hsu et al.~\cite{hsu2026noteevent} & fingering & notes & string, fret & -- & beats Fretting-Transformer \\
Chen et al.~\cite{chen2015technique}, TENT~\cite{su2019tent} & techniques & solo guitar & technique per note & no & TENT: +14.7\% note F \\
TabCNN~\cite{wiggins2019tabcnn} & end-to-end & solo guitar & frets per frame & no & GuitarSet baseline \\
FretNet~\cite{cwitkowitz2023fretnet} & end-to-end & solo guitar & string contours & no & finer pitch resolution \\
Noise2Fret~\cite{simionato2026noise2fret} & end-to-end & solo guitar & string, fret & no & beats baselines, GuitarSet/GOAT \\
TART~\cite{gupta2026tart} & hybrid & solo guitar & tab + techniques & no & 81.35 F50 / 71.8 / 54.08 Tab F1 \\
\bottomrule
\end{tabular}
\end{table*}

\subsection{Known and expected failure modes}
\emph{Distorted electric guitar.} Models trained on clean or classical guitar
lose accuracy on amplified tones~\cite{chen2022egdb}; in our measurements
(below) most of the added errors are notes estimated an octave too low.
Tone-diverse training data is the documented
remedy~\cite{chen2025egdbpg,pedroza2024tones}.

\emph{Separation residue.} Note models trained on solo guitar see bleed from
other instruments and smeared transients after separation, which become
false and missed onsets. No published guitar transcriber is trained on
separated stems.

\emph{Several guitars.} Separators produce one guitar stem and
multi-instrument transcribers one stream per instrument class; two guitars
playing different parts in the same region of the stereo image are merged.

\emph{Unisons and repeated pitches.} MT3-style tokenizers cannot represent the
same pitch twice on one instrument class at once~\cite{rouard2026muscriptor};
on guitar this occurs for unisons played on two strings and for overlapping
repeated notes.

\emph{Fingering conventions.} Symbolic fingering models learn the conventions
of their training tablatures; DadaGP consists of user-made tablatures, so
models pre-trained on it reproduce their choices, and they cannot use
timbre unless conditioned on audio.

\emph{Expressive and unusual rhythm.} Beat trackers fail on slow, expressive
or rubato music~\cite{ahn2026smc}; compound meters and tuplets are rarely
notated correctly by GTT systems.

\emph{Techniques under distortion.} Contour-based detection assumes a clean
pitch track and a visible pick attack; both degrade with distortion and
compression.

\subsection{Observations from our own measurements}
As a reference point for the use case, we scored the GAPS model and two
versions fine-tuned on electric guitar in many tones (EGDB, Guitar-TECHS and,
for the second, GuitarSet players 00--02) on held-out sets, with
\texttt{mir\_eval} F50 on raw notes. The fine-tuned model raised F50 on
GOAT amplifier renders, which it never saw, from 0.56 to 0.71 and halved
octave-low errors there (8\% to 5\% of notes), while staying at 0.89--0.92 on
GuitarSet. Fed band mixes (labelled guitar over MUSDB18 backing tracks)
without separation, every variant scored 0.50--0.51.

We then ran one available model per stage and approach on the same sets.
\emph{Separation:} transcribing the guitar stem raised the fine-tuned model
from 0.50 (raw mix) to 0.74 with \texttt{htdemucs\_6s} and 0.81 with a
BS-RoFormer six-stem checkpoint, close to its 0.83 on isolated amplified
guitar. \emph{Multi-instrument transcription:} YourMT3+ was the best model on
GuitarSet (F50 0.92--0.93, with much better offsets) but scored 0.39 on GOAT
amplifier renders and did not benefit from separation (0.51 on the mix, 0.52
on the BS-RoFormer stem). \emph{Beats:} on GuitarSet accompaniment, Beat This!
reached a beat F-measure of 0.97 with its larger checkpoint, but about 0.5 on
unaccompanied solo lines. \emph{Modular versus end-to-end:} with Tab F1 as the
metric, the modular pipeline (fine-tuned note model, separation, Viterbi
fingering) beat TabCNN and its tone-augmented variant on electric guitar
(0.55 vs 0.37 on EGDB amplifier tones; 0.59 vs 0.37 on separated mixes), while
TabCNN was better on GuitarSet solo lines (0.81 vs 0.53), where the Viterbi
placed only 56\% of the reference notes on the human's string. Noise2Fret's
released checkpoints reproduce its reported scale (tablature F 0.92 on
GuitarSet player 00) under its own window-based protocol, which scores only
windows containing an onset and is not comparable to note-level F1.

These numbers support three points of this review: tone augmentation
addresses distortion; band-split separation removes most of the mix penalty;
and fingering, not note detection, is the weakest stage on melodic lines.

\subsection{Open problems}
\begin{enumerate}
\item \emph{Mix-level tablature.} No method outputs string and fret from a
band recording. Extending a multi-instrument transcriber such as MuScriptor
with string and fret tokens would combine mix robustness with timbre-aware
fingering, and adding the string to each note token would also remove the
same-pitch limitation.
\item \emph{Separation for transcription.} Separators are evaluated by SDR,
not by the transcription accuracy of their stems; transcribers are not
trained on separated stems. Training each on the other's output is untested
for guitar.
\item \emph{Part separation.} Splitting several guitars of the same class
(rhythm and lead, doubled parts) needs either spatial cues or separators
conditioned on a part.
\item \emph{Techniques in context.} Learned technique classifiers trained on
GOAT and IDMT-SMT-Guitar with tone augmentation, evaluated on amplified and
polyphonic material.
\item \emph{A real benchmark.} A public set of commercial recordings with
aligned reference tablatures would make the use case measurable. MuScriptor's
alignment procedure (bar-line interpolation and DTW on chroma and onset
features with Sync Toolbox~\cite{muller2021sync}) shows a way to build one
from existing tablatures.
\item \emph{Rhythm notation.} Quantisation, tuplets and compound meters are
handled by rules in current systems and are not evaluated.
\end{enumerate}
